% Options for packages loaded elsewhere
\PassOptionsToPackage{unicode}{hyperref}
\PassOptionsToPackage{hyphens}{url}
\PassOptionsToPackage{dvipsnames,svgnames,x11names}{xcolor}
%
\documentclass[
  letterpaper,
  DIV=11,
  numbers=noendperiod]{scrartcl}

\usepackage{amsmath,amssymb}
\usepackage{iftex}
\ifPDFTeX
  \usepackage[T1]{fontenc}
  \usepackage[utf8]{inputenc}
  \usepackage{textcomp} % provide euro and other symbols
\else % if luatex or xetex
  \usepackage{unicode-math}
  \defaultfontfeatures{Scale=MatchLowercase}
  \defaultfontfeatures[\rmfamily]{Ligatures=TeX,Scale=1}
\fi
\usepackage{lmodern}
\ifPDFTeX\else  
    % xetex/luatex font selection
\fi
% Use upquote if available, for straight quotes in verbatim environments
\IfFileExists{upquote.sty}{\usepackage{upquote}}{}
\IfFileExists{microtype.sty}{% use microtype if available
  \usepackage[]{microtype}
  \UseMicrotypeSet[protrusion]{basicmath} % disable protrusion for tt fonts
}{}
\makeatletter
\@ifundefined{KOMAClassName}{% if non-KOMA class
  \IfFileExists{parskip.sty}{%
    \usepackage{parskip}
  }{% else
    \setlength{\parindent}{0pt}
    \setlength{\parskip}{6pt plus 2pt minus 1pt}}
}{% if KOMA class
  \KOMAoptions{parskip=half}}
\makeatother
\usepackage{xcolor}
\setlength{\emergencystretch}{3em} % prevent overfull lines
\setcounter{secnumdepth}{-\maxdimen} % remove section numbering
% Make \paragraph and \subparagraph free-standing
\ifx\paragraph\undefined\else
  \let\oldparagraph\paragraph
  \renewcommand{\paragraph}[1]{\oldparagraph{#1}\mbox{}}
\fi
\ifx\subparagraph\undefined\else
  \let\oldsubparagraph\subparagraph
  \renewcommand{\subparagraph}[1]{\oldsubparagraph{#1}\mbox{}}
\fi


\providecommand{\tightlist}{%
  \setlength{\itemsep}{0pt}\setlength{\parskip}{0pt}}\usepackage{longtable,booktabs,array}
\usepackage{calc} % for calculating minipage widths
% Correct order of tables after \paragraph or \subparagraph
\usepackage{etoolbox}
\makeatletter
\patchcmd\longtable{\par}{\if@noskipsec\mbox{}\fi\par}{}{}
\makeatother
% Allow footnotes in longtable head/foot
\IfFileExists{footnotehyper.sty}{\usepackage{footnotehyper}}{\usepackage{footnote}}
\makesavenoteenv{longtable}
\usepackage{graphicx}
\makeatletter
\def\maxwidth{\ifdim\Gin@nat@width>\linewidth\linewidth\else\Gin@nat@width\fi}
\def\maxheight{\ifdim\Gin@nat@height>\textheight\textheight\else\Gin@nat@height\fi}
\makeatother
% Scale images if necessary, so that they will not overflow the page
% margins by default, and it is still possible to overwrite the defaults
% using explicit options in \includegraphics[width, height, ...]{}
\setkeys{Gin}{width=\maxwidth,height=\maxheight,keepaspectratio}
% Set default figure placement to htbp
\makeatletter
\def\fps@figure{htbp}
\makeatother

\KOMAoption{captions}{tableheading}
\makeatletter
\makeatother
\makeatletter
\makeatother
\makeatletter
\@ifpackageloaded{caption}{}{\usepackage{caption}}
\AtBeginDocument{%
\ifdefined\contentsname
  \renewcommand*\contentsname{Table of contents}
\else
  \newcommand\contentsname{Table of contents}
\fi
\ifdefined\listfigurename
  \renewcommand*\listfigurename{List of Figures}
\else
  \newcommand\listfigurename{List of Figures}
\fi
\ifdefined\listtablename
  \renewcommand*\listtablename{List of Tables}
\else
  \newcommand\listtablename{List of Tables}
\fi
\ifdefined\figurename
  \renewcommand*\figurename{Figure}
\else
  \newcommand\figurename{Figure}
\fi
\ifdefined\tablename
  \renewcommand*\tablename{Table}
\else
  \newcommand\tablename{Table}
\fi
}
\@ifpackageloaded{float}{}{\usepackage{float}}
\floatstyle{ruled}
\@ifundefined{c@chapter}{\newfloat{codelisting}{h}{lop}}{\newfloat{codelisting}{h}{lop}[chapter]}
\floatname{codelisting}{Listing}
\newcommand*\listoflistings{\listof{codelisting}{List of Listings}}
\makeatother
\makeatletter
\@ifpackageloaded{caption}{}{\usepackage{caption}}
\@ifpackageloaded{subcaption}{}{\usepackage{subcaption}}
\makeatother
\makeatletter
\@ifpackageloaded{tcolorbox}{}{\usepackage[skins,breakable]{tcolorbox}}
\makeatother
\makeatletter
\@ifundefined{shadecolor}{\definecolor{shadecolor}{rgb}{.97, .97, .97}}
\makeatother
\makeatletter
\makeatother
\makeatletter
\makeatother
\ifLuaTeX
  \usepackage{selnolig}  % disable illegal ligatures
\fi
\IfFileExists{bookmark.sty}{\usepackage{bookmark}}{\usepackage{hyperref}}
\IfFileExists{xurl.sty}{\usepackage{xurl}}{} % add URL line breaks if available
\urlstyle{same} % disable monospaced font for URLs
\hypersetup{
  pdftitle={One-dimensional Turbulence Model (notes)},
  pdfauthor={Sienna White},
  colorlinks=true,
  linkcolor={blue},
  filecolor={Maroon},
  citecolor={Blue},
  urlcolor={Blue},
  pdfcreator={LaTeX via pandoc}}

\title{One-dimensional Turbulence Model (notes)}
\author{Sienna White}
\date{}

\begin{document}
\maketitle
\ifdefined\Shaded\renewenvironment{Shaded}{\begin{tcolorbox}[sharp corners, breakable, enhanced, frame hidden, borderline west={3pt}{0pt}{shadecolor}, interior hidden, boxrule=0pt]}{\end{tcolorbox}}\fi

\section{Problem 3}

Write the equation of motion for a velocity fluctuation.

We start by recalling the equation of motion for momentum, after
performing Reynold's decomposition and splitting components into their
mean and fluctuating parts\ldots{}

\[ \partial_t (\bar{u_j} + u_j') + (\bar{u_i} + u_i') \partial_i (\bar{u_j} + u_j') = - \frac{1}{\rho} \partial_j (\bar{P} +P') + \nu \partial_k \partial_k (\bar{u_j} + u_j')\]

\begin{equation}
     \partial_t \bar{u_j}  + \partial_t u_j' + \bar{u_i} \partial_i \bar{u_j} + \bar{u_i} \partial_i u_j' + u_i' \partial_i \bar{u_j}  + u_i' \partial_i u_j' = 
     - \frac{1}{\rho} \partial_j \bar{P} - \frac{1}{\rho} \partial_j P' + \nu \partial_k \partial_k \bar{u_j} +\nu \partial_k \partial_k u_j'
\end{equation}

Now we can subtract our RANS equation for momentum from this equation to
isolate the turbulent components (\(u'\)). Recall our RANS equation for
momentum reads

\begin{equation}
     \partial_t \bar{u_j}  + \bar{u_i} \partial_i \bar{u_j}  =  - \frac{1}{\rho} \partial_j \bar{P}  + \nu \partial_k \partial_k \bar{u_j} -  \partial_i \overline{u_j'u_i'}  . 
\end{equation}

Subtracting Equation (2) from Equation (1) yields \begin{equation}
     \partial_t u_j' + \bar{u_i} \partial_i u_j' + u_i' \partial_i \bar{u_j}  + u_i' \partial_i u_j' =  - \frac{1}{\rho} \partial_j P'  +\nu \partial_k \partial_k u_j' + \partial_i \overline{u_j'u_i'}  . 
\end{equation}

\begin{equation}
     \partial_t u_j' + \bar{u_i} \partial_i u_j'  =  - \frac{1}{\rho} \partial_j P'  +\nu \partial_k \partial_k u_j' + \partial_i \overline{u_j'u_i'} - u_i' \partial_i \bar{u_j}  - u_i' \partial_i u_j'  . 
\end{equation}

\subsection{Derive equation for Reynold's Stress Transport}

We multiply this equation by \(u_k'\), expanding this into a 3x3 tensor
equation. \begin{equation}
     u_k' \partial_t u_j' + u_k' \bar{u_i} \partial_i u_j'  =  - u_k' \frac{1}{\rho} \partial_j P'  + u_k' \nu \partial_z \partial_z u_j' + u_k' \partial_i \overline{u_j'u_i'} - u_k' u_i' \partial_i \bar{u_j}  - u_k' u_i' \partial_i u_j'  . 
\end{equation}

Reynold's Average: \begin{equation}
     \langle u_k' \partial_t u_j' \rangle  +
     \langle u_k' \bar{u_i} \partial_i u_j' \rangle  =  
      - \langle u_k' \frac{1}{\rho} \partial_j P' \rangle + 
     \langle u_k' \nu \partial_z \partial_z u_j' \rangle  + 
     \langle u_k' \partial_i \overline{u_j'u_i'} \rangle - 
     \langle u_k' u_i' \partial_i \bar{u_j} \rangle - 
     \langle u_k' u_i' \partial_i u_j' \rangle . 
\end{equation} The divergence of the averaged Reynold's Stress goes to
zero, as it's an averaged quantity multiplied by a fluctuating value

\begin{equation}
     \langle u_k' \partial_t u_j' \rangle  +
     \langle u_k' \bar{u_i} \partial_i u_j' \rangle  =  
      - \langle u_k' \frac{1}{\rho} \partial_j P' \rangle + 
     \langle u_k' \nu \partial_z \partial_z u_j' \rangle  + 
     \langle u_k' u_i' \partial_i \bar{u_j} \rangle - 
     \langle u_k' u_i' \partial_i u_j' \rangle . 
\end{equation}

We now take a second equation, flip the indices \begin{equation}
     \langle u_j' \partial_t u_k' \rangle  +
     \langle u_j' \bar{u_i} \partial_i u_k' \rangle  =  
      - \langle u_j' \frac{1}{\rho} \partial_k P' \rangle + 
     \langle u_j' \nu \partial_z \partial_z u_k' \rangle  + 
     \langle u_j' u_i' \partial_i \bar{u_k} \rangle - 
     \langle u_j' u_i' \partial_i u_k' \rangle 
\end{equation}

and add these together. To reduce the messiness, we will first look at
the LHS,

\begin{equation}
     \langle u_j' \partial_t u_k' \rangle  +
     \langle u_j' \bar{u_i} \partial_i u_k' \rangle  +
     \langle u_k' \partial_t u_j' \rangle  +
     \langle u_k' \bar{u_i} \partial_i u_j' \rangle.
\end{equation}

To simply, we apply the reverse product rule (ie, given we know)
\[ \frac{\partial}{\partial t} \langle u_k' u_j' \rangle               = \langle u_k' \frac{\partial}{\partial t}  u_j' \rangle + 
    \langle u_j' \frac{\partial}{\partial t}  u_k' \rangle \]

and
\[ \bar{u_i}\frac{\partial}{\partial x_i} \langle u_k' u_j' \rangle  = \bar{u_i} \langle u_k' \frac{\partial}{\partial x_i} u_j' \rangle  + \bar{u_i} \langle  u_j' \frac{\partial}{\partial x_i} u_k' \rangle,
\]

we can rewrite this LHS as

\[ \partial_t \langle u_k' u_j' \rangle  + \bar{u_i} \partial_i \langle u_k' u_j' \rangle. \]

Now onto the fun bit\ldots{} the RHS. Currently we have

\[ - \langle u_k' \frac{1}{\rho} \partial_j P' \rangle  
     - \langle u_j' \frac{1}{\rho} \partial_k P' \rangle  
    + \langle u_k' \nu \partial_z \partial_z u_j' \rangle  
   + \langle u_j' \nu \partial_z \partial_z u_k' \rangle  
    +   \langle u_k' u_i' \partial_i \bar{u_j} \rangle 
    + \langle u_j' u_i' \partial_i \bar{u_k} \rangle  
   -  \langle u_k' u_i' \partial_i u_j' \rangle  
  - \langle u_j' u_i' \partial_i u_k' \rangle . 
\]

Some of this follows the same pattern (reverse product rule), allowing
us to expand our fluctuating pressure term. Consider that
\[ \frac{\partial }{\partial x_j} \langle u_k' P' \rangle =
\langle u_k' \partial_j P' \rangle  + \langle P' \partial_j u_k' \rangle.\]

We substitute \[ - \langle u_k' \frac{1}{\rho} \partial_j P' \rangle = 
 \frac{1}{\rho} \left( \langle P' \partial_j u_k' \rangle - \partial_j \langle u_k' P' \rangle \right) \]

\[ - \langle u_j' \frac{1}{\rho} \partial_k P' \rangle = 
 \frac{1}{\rho} \left( \langle P' \partial_k u_j' \rangle - \partial_k \langle u_j' P' \rangle \right) \]

Next we identify the advection of mean flow by the Reynolds stress,
i.e., the terms

\[    \langle u_k' u_i' \partial_i \bar{u_j} \rangle 
    + \langle u_j' u_i' \partial_i \bar{u_k} \rangle. \]

We won't rearrange these terms for now, since they're already in a form
involving an averaged Reynold's stress. The viscous dissipation,
however, looks like an interesting challenge,

\[ \nu  \langle u_k' \partial_z \partial_z u_j' \rangle  
   +  \nu \langle u_j' \partial_z \partial_z u_k' \rangle. \]

We would like to write this as the diffusion of the Reynolds stress. We
expect such diffusion to look like
\[ \nu \partial_z \partial_z \langle  u_k'  u_j' \rangle ,\]

which would expand (ignoring the coefficient) as

\[ \partial_z \left( \partial_z \langle  u_k'  u_j' \rangle \right) = \partial_z  ( \langle u_k' \partial_z  u_j' \rangle + \langle u_j' \partial_z u_k' \rangle) \]

\[ =  \langle u_k' \partial_z \partial_z u_j' \rangle + \langle \partial_z u_j'\partial_z u_k' \rangle +  \langle \partial_z u_j'\partial_z u_k' \rangle + \langle u_j' \partial_z  \partial_z u_k' \rangle.\]

We relate this to the term in the equation by substituting

\[ \nu  \langle u_k' \partial_z \partial_z u_j' \rangle  
+  \nu \langle u_j' \partial_z \partial_z u_k' \rangle =  \nu \partial_z \partial_z \langle  u_k'  u_j' \rangle - 2 \nu \langle \partial_z u_j'\partial_z u_k' \rangle.
\]

The last terms we want to deal with are the advection of turbulent
fluctuations by the Reynolds stress itself,
\[  -  \langle u_k' u_i' \partial_i u_j' \rangle  
  - \langle u_j' u_i' \partial_i u_k' \rangle . \]

We recognize this pattern, and can apply reverse-product-rule. The last
step is due to impressibility (allows us to put \(u_i\) inside the
derivative).\\
\[  -  \langle u_i' u_k'  \partial_i u_j' \rangle  
  - \langle u_i' u_j' \partial_i u_k' \rangle =  -  \langle u_i' \partial_i u_k' u_j' \rangle  = -  \langle  \partial_i u_k' u_j' u_i' \rangle \]

Now we can put together all the steps. For notation in this last step, I
am using capital \(U\) to denote an averaged velocity (instead of an
overbar).

\begin{equation}
    \partial_t \langle u_k' u_j' \rangle  + U_i \partial_i \langle u_k' u_j' \rangle= 
\end{equation}

\begin{verbatim}
% triple correlation 
\end{verbatim}

\[  -   \partial_i \langle  u_k' u_j' u_i' \rangle 
    % transport mean flow by RE 
    +  \langle u_k' u_i'\rangle  \partial_i U_j  
    + \langle u_j' u_i' \rangle \partial_i U_k  
    \]

\begin{verbatim}
 % viscous stress
\end{verbatim}

\[  + \nu \partial_z \partial_z \langle  u_k'  u_j' \rangle - 2 \nu \langle \partial_z u_j'\partial_z u_k' \rangle +
    % pressure 
\frac{1}{\rho} \left( P' \langle \partial_j u_k' + \partial_k u_j' \rangle - \partial_j \langle u_k' P' \rangle  - \partial_k \langle u_j' P' \rangle \right)  \]

We can define all these terms and define them from a physics
perspective. Starting left to right,

\begin{enumerate}
    \item Storage / time variation of the Reynold's stress
    \item Advection of the Reynold's stress by the mean flow
    \item Triple Correlation: transport of Reynold's stress by turbulent motions. 
    \item Advection of flow by Reynold's stress. This is also seen as a production term, as it captures the conversion of turbulence from the mean flow through shear. 
    \item (shear production)
    \item Diffusion of mometum from viscosity / molecular interactions. 
    \item This is a very exciting term! It is the viscous dissipation term and it \textit{directly destroys} Reynold's stress through conversion to heat.  
    \item These next two terms are very interesting. Franko spent a long time talking to me about this one before he graduated, and said if I understood this term, I would ``understand some grand part of turbulence." Whether or not that is true has yet to be seen. However, this is the pressure-strain correlation. It resembles the rate-of-strain tensor, and similar to $S_{ij},$ this tensor is traceless. This term cannot create or destroy turbulence, however, it often performs what is called ``pressure scrambling." One of these terms will anisotropize the turbulence (``fast term") and the other isotropizes the flow (``slow term"). 
    \item These terms are transport by a pressure-fluctuation correlation term. 
\end{enumerate}

The TKE Equation can be expressed as \begin{equation}
\frac{\partial k}{\partial t} + \bar{U_j}
\end{equation}

The density profile is calculated according to \begin{equation}
\rho = \rho_0 (1 - \alpha (C - T_0))
\end{equation}

\section{The Mellor-Yamada closure}

We start with the continuity equation, \begin{equation}
\frac{\partial \rho}{\partial t} + \frac{\partial }{\partial x_i} (\rho U_i) = 0 
\end{equation} where flow is considered to be divergent-free, or
solenoidal.

Meanwhile, the Reynolds-Averaged Navier Stokes reads: \begin{equation}
\frac{\partial \overline{U}}{\partial t} 
+ \overline{U_j} \frac{\partial \overline{U_k} }{\partial x_j} 
+ \rho \epsilon_{jkl} f_k U_l 
= \frac{\partial }{\partial x_k} \overline{u_k u_j} 
- \frac{\partial \overline{P}}{\partial x_j} 
- g \hat{e_3} 
\end{equation}

The Boussinesq approximation is invoked to represent that density
fluctuaions affect bouyancy but no other transport/state terms. These
fluctuations are written \(\beta \theta\) , where \(\beta\) is the
coefficient of thermal expansion.

Mellor and Herring (1973) uses the following equation: \begin{equation}
\langle \frac{p}{\rho} (\frac{\partial u_i}{\partial x_j} + \frac{\partial u_j}{\partial x_i})  \rangle = - \frac{q}{3l_1} ( \overline{u_iu_j} - \frac{\delta_{ij}}{3} q^2) + C_1 q^2 (\frac{\partial U_i}{\partial x_j} + \frac{\partial U_j}{\partial x_i})  
\end{equation}

We are establishing equations for the following:
\(\overline{u_k u_i u_j},~ \overline{p u_j},~ \overline{u_i u_j \theta},\)
and \(\overline{p\theta}\)

Math Note The velocity gradient \(\frac{\partial U_j}{\partial x_i}\)
can be decomposed into symmetric and antisymmetric parths,
\[\frac{\partial U_j}{\partial x_i} = S_{ij} + \Omega_{ij} \] where the
symmetric rate-of-strain tensor is
\[ S_{ij} = \frac{1}{2} \left( \frac{\partial U_i}{\partial x_j} + \frac{\partial U_j}{\partial x_i} \right) \]\\
and the asymmetric part is,
\[ \Omega_{ij} = \frac{1}{2} \left( \frac{\partial U_i}{\partial x_j} - \frac{\partial U_j}{\partial x_i} \right) \]

\subsection{Parameterization}

\begin{equation}
\langle u_k u_i u_j \rangle = \frac{3}{5} 
 ~l q S_q 
 \left( \frac{\partial \langle u_i u_j \rangle}{\partial x_k}  
+ \frac{\partial \langle u_i u_k \rangle}{\partial x_j}
+ \frac{\partial \langle u_j u_k \rangle}{\partial x_i} \right)
\end{equation}

In this equation \(S_q\) is a dimensionless number.

Sm : \begin{equation}
\frac{A_1}{A_2} \frac{B_1(\gamma - C_1) - (B_1 (\gamma_1 - C_1) + 6 (A_1 + 3A_2)) R_f}{B_1 \gamma_1 - (B_1(\gamma_1 + \gamma_2) - 3A_1) R_f}
\end{equation}

where \[ \gamma_1 = \frac{1}{3} - \frac{2 A_1}{B_1} \] and
\[ \gamma_2 = \frac{B_2}{B_1} - \frac{6 A_1}{B_1} \]

\[ \]

Sh:
\[ S_h  =  \frac{A_2 (1 - \frac{6 A_1}{B_1})}{1-(3 A_2 B_2 G_h + 18 A_1 A_2 G_h) }  = \frac{A_2 (1 - \frac{6 A_1}{B_1})}{(1-3 A_2 G_h (B_2+6 A_1))}   \]

\[ 
S_m = \frac{A_1 (1 - 3C_1 - \frac{6A_1}{B_1}) + S_H G_h(18A_1 + 9A_1A_2)}{(1 - 9A_1 A_2 G_h)}
 \]

Richardson's number:
\[ G_H = - \frac{L^2}{Q^2}\frac{g}{\rho_0}\frac{\partial}{\partial z}\]

Turbulent visosity: \begin{equation}
K_q = S_m  Q  L  \nu
\end{equation}

Turbulent diffusivity: \begin{equation}
K_q = S_q  Q  L  \nu 
\end{equation}

Turbulent scalar diffusivity: \begin{equation}
K_q = S_h * Q * L * \nu 
\end{equation}

\subsection{Chosen parameters}

\[ C_1 = \gamma_1 - \frac{1}{3A_1} B_1^{1/3} \approx 0.08\]

\[ B_2 = \frac{B_1^{1/3}}{P_{rt}} \frac{u_T^2 \langle \theta^2 \rangle}{H^2} \approx \]
\[ B_1 = 16.6  \]

\section{writing equations based on code}

\[ bX = bX - \frac{aX cX}{bX}\]

\[ dX = dX - \frac{aX dX}{bX}\]

\[ u_* = U_{bed} \sqrt{C_d}\]

\[G_h = - \left(\frac{N_{BV} * L}{Q + Small} \right)^2\]

\[ diss= 2 \frac{dt * \sqrt{Q^2}}{B_1 L_p} \]

\begin{tikzpicture}
  % Matrix
  \matrix (m) [matrix of math nodes, nodes in empty cells,
    row sep=0.5em, column sep=0.5em,
    nodes={minimum size=1.5em, anchor=center, outer sep=0}]
  {
    1 & 2 & 3 & 4 \\
    5 & 6 & 7 & 8 \\
    9 & 10 & 11 & 12 \\
    13 & 14 & 15 & 16 \\
  };

  % Highlighting diagonals
  \draw[red, thick] (m-1-1.north west) -- (m-4-4.south east); % Main diagonal
  \draw[blue, thick] (m-1-2.north west) -- (m-4-3.south east); % Upper diagonal
  \draw[green, thick] (m-2-1.north west) -- (m-3-4.south east); % Lower diagonal
\end{tikzpicture}



\end{document}
